% called by main.tex
%
\chapter{Aspectos éticos y limitaciones}
\label{ch::etica}

Este capítulo recoge las consideraciones éticas que han guiado el trabajo y una revisión
explícita de sus limitaciones, que delimitan el alcance de cualquier conclusión.

\section{Consideraciones éticas}

El sistema estudia decisiones económicas automatizadas y es especialmente sensible al
sobreajuste cuando el horizonte de datos es reducido. La actuación responsable es no operar
con capital real y no convertir un diagnóstico \emph{post-hoc} en una afirmación de
rentabilidad. La política completa queda congelada para observación prospectiva, sin que los
días usados para formularla cuenten de nuevo como validación.

En cuanto a los datos, se adopta un marco de confidencialidad, integridad y disponibilidad.
Una referencia externa retrasada o errónea puede inducir una señal espuria; por ello el
sistema se entrena solo con sesiones aceptadas, separa datos crudos, características y
etiquetas, mantiene trazabilidad de los artefactos e inhibe la operativa ante datos
obsoletos. El trabajo emplea datos públicos de mercado y no datos personales. El corpus
completo no se distribuye por tamaño; se publican el esquema, un ledger anonimizado de
acciones y el código de evaluación.

\section{Limitaciones}

\textbf{Selección posterior al bloque.} El especialista base sí estaba congelado y su
resultado limpio es $+0{,}349$ ticks por acción al coste de referencia. El filtro de
volatilidad se incorporó después de observar el cambio de régimen; por tanto, su resultado de
$+1{,}069$ ticks sobre el mismo bloque es exploratorio y necesita fechas posteriores.

\textbf{Soporte temporal.} El conjunto fuera de muestra abarca cinco días. Tres de cinco son
positivos para la política base y cinco de cinco dentro del subconjunto diagnóstico. Ninguna
de las dos cifras constituye una validación definitiva.

\textbf{Dependencia y solapamiento.} Las 318 unidades del diagnóstico pertenecen a 156
mercados, con hasta ocho señales por mercado y horizontes H60 solapados. El
\emph{bootstrap} por filas subestima esa dependencia. Un remuestreo agrupado por mercado
amplía el IC90\,\% a $[+0{,}193,+2{,}004]$, pero sigue describiendo un subconjunto elegido
\emph{post-hoc}. Además, no se ha fijado una regla de exposición simultánea; por ello la suma
acumulada no es PnL de cartera.

\textbf{Validación \emph{offline}.} La evaluación usa un modelo de ejecución conservador pero
simplificado. No reconstruye la cola orden por orden ni incorpora rellenos parciales. El
registro de sombra actual es retrospectivo; falta una ejecución prospectiva de mayor
fidelidad.

\textbf{Reproducibilidad parcial.} El repositorio permite recalcular las tablas públicas desde
el ledger anonimizado y probar la lógica compartida. La reproducción integral del
entrenamiento requiere el corpus privado de gran tamaño y se documenta como limitación, no
como capacidad publicada.

\textbf{Patrón en U.} Dentro de la baja volatilidad se observa una posible relación en forma de
U entre volatilidad y resultado. Como se descubrió en los mismos cinco días, se conserva solo
como hipótesis para trabajo futuro.
